\documentclass[10pt,a4paper]{amsart}
\usepackage[margin=19mm]{geometry}
\usepackage{amsmath,amssymb,mathtools}
\usepackage[scr=rsfs,cal=euler]{mathalfa}
\usepackage{xfrac}
\usepackage[unicode,hidelinks,bookmarks=false]{hyperref}
\newcommand{\ie}{i.e.\ }
\newcommand{\cf}{cf.\ }
\newcommand{\resp}{respectively\ }
\newcommand{\Subsection}{\subsection}
\providecommand{\nobreakdash}{\nobreak\hbox{-}}
\setlength{\emergencystretch}{2em}
\multlinegap0pt
\usepackage{bm}
\usepackage{bbm}
\usepackage{dsfont}
\usepackage{xspace}
\usepackage{cancel}
\usepackage{mathtools}
\usepackage{amsthm}
\makeatletter
\@ifundefined{theorem}{\newtheorem{theorem}{Theorem}}{}
\makeatother
\title[Curvature Preserving Fractal Interpolation Functions: A Hybr]{Curvature Preserving Fractal Interpolation Functions: A Hybrid Geometric Approach}
\author{K R Tyada}
\date{Source: arXiv:2602.01707v1 (CC BY 4.0)}
\begin{document}
\maketitle

\noindent\emph{Source and licence.} Curvature Preserving Fractal Interpolation Functions: A Hybrid Geometric Approach. Version arXiv:2602.01707v1 (CC BY 4.0), distributed under the Creative Commons Attribution 4.0 International licence, as recorded in the arXiv licence metadata for this exact version and shown on the abstract page https://arxiv.org/abs/2602.01707v1. Full rights notice: licenses/arxiv-2602.01707-CC-BY-4.0.txt.\par\smallskip

\noindent\textit{Editorial scope.} Counted unit: the Theorem whose source label is \texttt{None} in the article (source0.tex lines 237--246, source label \texttt{None}), delivered together with the article's own complete original proof of it (source0.tex lines 248--278, one \texttt{proof} environment of 31 lines). No mathematics is written, repaired or re-derived: statement and proof are the article's own text line for line. Required context retained uncounted: none. Every definition, display and label of those lines is retained unchanged. Omitted from this package and not redistributed: 1--236, 247--247, 279--422. Nothing inside a retained excerpt is omitted or altered: every retained body is delivered exactly as the article's lines are written, so no omission receipt of any kind accompanies this package. Citations: the retained excerpts carry no citation command at all, so no bibliography entry of the article is transcribed and no reference is redistributed. Reference adaptation: none was needed and none was made; every reference inside the delivered unit resolves inside it, so no unresolved reference or citation is produced. Printed slips kept literal: no printed word, symbol, hypothesis or number of the counted statement or of its proof was altered. Layout and class: the article's own class is \texttt{llncs} with packages including fontenc, amsmath, amssymb, epsfig, subfig, graphicx; the wrapper is the lane's pinned \texttt{amsart} editorial wrapper at 10pt on A4, which restates the theorem-like environments the retained text needs and adds the article's own macro definitions that the retained text needs (none), with extra wrapper packages (bm, bbm, dsfont, xspace, cancel, mathtools). The delivered numbering is the unit's own. Assets: no retained span contains a figure environment, a graphics-inclusion command, a PDF-page inclusion, a table or a TikZ picture; no image file is copied into this package, the package declares no asset dependency and no image file is redistributed with it. Licence: version arXiv:2602.01707v1 of this article is distributed under the Creative Commons Attribution 4.0 International licence, as recorded in the arXiv licence metadata for this exact version and shown on the abstract page https://arxiv.org/abs/2602.01707v1. A full rights notice is delivered at licenses/arxiv-2602.01707-CC-BY-4.0.txt. Characters: every retained excerpt is pure ASCII, so the delivered engine, XeLaTeX with its default Latin Modern text font, reports no missing character.
	\begin{theorem}[Curvature Stability]
		Let \( F \) and \( \bar{F} \) be two cubic fractal interpolation functions defined on the interval \( I = [x_1, x_N] \), generated by the IFS mappings:
		\[
		F(L_s(x)) = \lambda_s F(x) + P_s(x), \quad \bar{F}(L_s(x)) = \bar{\lambda}_s \bar{F}(x) + P_s(x),
		\]
		for \( s \in \Im^* \), where \( |\lambda_s|, |\bar{\lambda}_s| < 1 \), and \( P_s(x) \) is the cubic polynomial satisfying the endpoint conditions. Let \( \kappa_F \) and \( \kappa_{\bar{F}} \) denote the curvatures of \( F \) and \( \bar{F} \), respectively. Then, for sufficiently small perturbations \( \|\bar{\lambda} - \lambda\|_\infty = \max_s |\bar{\lambda}_s - \lambda_s| \), there exist constants \( C_1, C_2 > 0 \) such that:
		\[
		|\kappa_F(x) - \kappa_{\bar{F}}(x)| \leq C_1 \|\bar{\lambda} - \lambda\|_\infty + C_2 \|\bar{\lambda} - \lambda\|_\infty^2, \quad \forall x \in I.
		\]
	\end{theorem}

	\begin{proof}
		Let \( e(x) = F(x) - \bar{F}(x) \) denote the difference between the two interpolants. From the IFS equations, the value of \( e \) at a subinterval satisfies:
		\[
		e(L_s(x)) = F(L_s(x)) - \bar{F}(L_s(x)) = \lambda_s F(x) + P_s(x) - \left( \bar{\lambda}_s \bar{F}(x) + P_s(x) \right).
		\]
		Simplifying and adding and subtracting \( \lambda_s \bar{F}(x) \) yields:
		\[
		e(L_s(x)) = \lambda_s (F(x) - \bar{F}(x)) + (\lambda_s - \bar{\lambda}_s) \bar{F}(x) = \lambda_s e(x) - \delta\lambda_s \cdot \bar{F}(x),
		\]
		where \( \delta\lambda_s = \bar{\lambda}_s - \lambda_s \). This functional equation for \( e(x) \) is itself a fractal function. Since \( \max_s |\lambda_s| \leq \alpha < 1 \) and \( \bar{F} \) is bounded on the compact interval \( I \), it follows from the properties of iterated function systems that:
		\[
		\|e\|_\infty = \mathcal{O}(\|\bar{\lambda} - \lambda\|_\infty).
		\]
		By differentiating the IFS relation and applying similar reasoning to \( e'(x) \) and \( e''(x) \), we obtain the corresponding bounds for the derivatives:
		\[
		\|e'\|_\infty = \mathcal{O}(\|\bar{\lambda} - \lambda\|_\infty), \quad \|e''\|_\infty = \mathcal{O}(\|\bar{\lambda} - \lambda\|_\infty).
		\]
		
		The curvature is given by \( \kappa_F(x) = \frac{F''(x)}{(1 + (F'(x))^2)^{3/2}} \). Consider the difference:
		\[
		|\kappa_F(x) - \kappa_{\bar{F}}(x)| = \left| \frac{F''}{(1 + (F')^2)^{3/2}} - \frac{\bar{F}''}{(1 + (\bar{F}')^2)^{3/2}} \right|.
		\]
		Adding and subtracting \( \frac{\bar{F}''}{(1 + (F')^2)^{3/2}} \) allows us to split the difference:
		\[
		|\kappa_F - \kappa_{\bar{F}}| \leq \underbrace{\frac{|F'' - \bar{F}''|}{(1 + (F')^2)^{3/2}}}_{T_1} + \underbrace{|\bar{F}''| \left| \frac{1}{(1 + (F')^2)^{3/2}} - \frac{1}{(1 + (\bar{F}')^2)^{3/2}} \right|}_{T_2}.
		\]
		Since \( F' \) is bounded away from infinity, the denominator in \( T_1 \) is bounded below, implying \( T_1 = \mathcal{O}(\|e''\|_\infty) = \mathcal{O}(\|\bar{\lambda} - \lambda\|_\infty) \). For \( T_2 \), the function \( \phi(a) = (1+a^2)^{-3/2} \) is continuously differentiable, and \( \bar{F}'' \) is bounded. By the Mean Value Theorem, \( T_2 = \mathcal{O}(|F' - \bar{F}'|) = \mathcal{O}(\|e'\|_\infty) = \mathcal{O}(\|\bar{\lambda} - \lambda\|_\infty) \). The potential quadratic term arises from a more detailed expansion of \( T_2 \) when considering the interaction between the errors in the function and its derivatives. Combining these estimates yields the desired result.
		%		\[
		%		|\kappa_F(x) - \kappa_{\bar{F}}(x)| \leq C_1 \|\bar{\lambda} - \lambda\|_\infty + C_2 \|\bar{\lambda} - \lambda\|_\infty^2.
		%		\]
	\end{proof}

\end{document}
